% ============================================================================
% PLANTILLA 09 — CASO DE ESTUDIO / ANÁLISIS DE DATOS
% ----------------------------------------------------------------------------
% Para: análisis de casos, interpretación de datos, gráficos y tablas
% estadísticas. Incluye figura pgfplots y tabla de datos con siunitx.
% Compilar: latexmk -pdf plantilla-09-caso-estudio.tex
% Solucionario: \documentclass[soluciones]{academic-exam}
% ============================================================================
\documentclass{academic-exam}
\usepackage{pgfplots}
\pgfplotsset{compat=1.17}

\universidad{Universidad Nacional de San Cristóbal de Huamanga}
\facultad{Facultad de Ciencias Económicas, Administrativas y Contables}
\escuela{Escuela Profesional de Economía}
\curso{Política Económica}
\codigocurso{EC-442}
\docente{Mg. Nombre Apellido}
\periodo{Semestre 2026-I}
\tipoevaluacion{Evaluación de Caso 01}
\fechaevaluacion{21 de julio de 2026}
\duracion{110 minutos}
\modalidad{Presencial · grupal}

\begin{document}
\membrete

\begin{instrucciones}
\begin{itemize}
  \item Lea íntegramente el caso antes de responder. Todas las preguntas
        se refieren al material presentado (texto, tabla y figura).
  \item Puede consultar apuntes de clase; no se permite internet.
  \item Fundamente cada respuesta con los datos del caso: las opiniones
        sin evidencia no reciben puntaje.
\end{itemize}
\end{instrucciones}

% ================================================================== CASO ===
\section*{\color{colorprimario}El caso · Choque inflacionario y respuesta
de política, 2021--2024}

Entre 2021 y 2023 la economía peruana experimentó el episodio
inflacionario más intenso en dos décadas. El alza de los precios
internacionales de alimentos y combustibles tras la pandemia y la guerra
en Ucrania, sumada a factores internos, llevó la inflación interanual de
Lima Metropolitana desde niveles cercanos a 2\,\% hasta un máximo de
8{,}8\,\% en junio de 2022, muy por encima del rango meta del BCRP
(1\,\%--3\,\%). El banco central respondió con el ciclo de alzas más
agresivo de su historia: la tasa de referencia pasó de 0{,}25\,\% en
julio de 2021 a 7{,}75\,\% en enero de 2023, donde permaneció hasta
iniciar el ciclo de recortes en setiembre de 2023.

\begin{datos}[Cuadro 1 · Indicadores macroeconómicos seleccionados]
\centering
\small
\begin{tabular}{@{}l S[table-format=2.1] S[table-format=2.2]
                 S[table-format=2.1] S[table-format=2.1]@{}}
  \toprule
  \textbf{Año} & {\textbf{Inflación} (\%)} & {\textbf{Tasa ref.} (\%, dic.)}
   & {\textbf{PBI} (var.\ \%)} & {\textbf{Expectativas} (\%)}\\
  \midrule
  2020 &  2.0 & 0.25 & -10.9 & 1.7\\
  2021 &  6.4 & 2.50 &  13.4 & 3.7\\
  2022 &  8.5 & 7.50 &   2.7 & 4.3\\
  2023 &  3.2 & 6.75 &  -0.6 & 3.1\\
  2024 &  2.0 & 5.00 &   3.3 & 2.5\\
  \bottomrule
\end{tabular}
\par\smallskip
{\footnotesize Fuente: elaborado con datos del BCRP (cifras referenciales
con fines pedagógicos). Expectativas de inflación a 12 meses, diciembre.}
\end{datos}

\begin{figure}[h]
\centering
\begin{tikzpicture}
\begin{axis}[
  width=0.86\textwidth, height=6.2cm,
  xlabel={Año}, ylabel={Porcentaje},
  xtick={2020,2021,2022,2023,2024},
  xticklabel style={/pgf/number format/1000 sep=},
  ymin=0, ymax=10,
  legend style={draw=none, fill=none, at={(0.98,0.95)}, anchor=north east,
                font=\small},
  grid=major, grid style={gray!20},
  every axis plot/.append style={line width=1pt, mark size=2pt},
]
\addplot[color=colorenfasis, mark=*]
  coordinates {(2020,2.0) (2021,6.4) (2022,8.5) (2023,3.2) (2024,2.0)};
\addlegendentry{Inflación interanual}
\addplot[color=colorprimario, mark=square*]
  coordinates {(2020,0.25) (2021,2.50) (2022,7.50) (2023,6.75) (2024,5.00)};
\addlegendentry{Tasa de referencia (dic.)}
\addplot[color=gristexto, dashed, no marks]
  coordinates {(2020,3) (2024,3)};
\addlegendentry{Techo del rango meta}
\end{axis}
\end{tikzpicture}
\caption{Inflación y tasa de política monetaria, 2020--2024.}
\end{figure}

% ============================================================ PREGUNTAS ====
\begin{pregunta}[5][Interpretación de datos]
Con el Cuadro~1 y la Figura~1, describa la cronología del episodio:
¿cuándo se desancla la inflación del rango meta, cuándo reacciona la
política monetaria y con qué rezago retorna la inflación al rango?
Calcule la tasa de interés \emph{real ex ante} de diciembre de 2022
(use las expectativas del cuadro) e interprete su signo.
\lineasrespuesta{10}
\begin{solucion}
La inflación supera el techo de 3\,\% durante 2021 (cierra en 6,4\,\%);
el BCRP inicia las alzas en agosto-2021 y alcanza el máximo de la tasa a
inicios de 2023; la inflación retorna al rango recién en 2024: un rezago
de transmisión de 18--24 meses, consistente con la literatura. Tasa real
ex ante dic-2022: $i - \pi^{e} = 7{,}50 - 4{,}3 = 3{,}2\,\%$, claramente
positiva y por encima de la tasa neutral estimada ($\approx 1{,}5\,\%$):
la postura era contractiva.
\end{solucion}
\end{pregunta}

\begin{pregunta}[5][Diagnóstico]
¿El choque inflacionario fue de oferta, de demanda o mixto? Sustente con
al menos dos hechos del caso y explique por qué el diagnóstico importa
para la respuesta de política.
\lineasrespuesta{10}
\begin{solucion}
Predominó el choque de oferta importado (alimentos y energía; la
inflación sube en 2021--2022 mientras el PBI se desacelera de 13,4\,\% a
2,7\,\%), con un componente de demanda por la recuperación pospandemia y
los retiros de fondos. El diagnóstico importa porque ante choques de
oferta la política monetaria no puede evitar el alza transitoria del
nivel de precios: su papel es impedir efectos de segunda vuelta anclando
expectativas — lo que justifica la reacción cuando las expectativas
(4,3\,\% en 2022) superaron el rango meta.
\end{solucion}
\end{pregunta}

\begin{pregunta}[5][Evaluación de política]
Un congresista sostuvo que el BCRP \enquote{frenó la economía
innecesariamente, pues la inflación igual iba a bajar cuando se
normalizaran los precios internacionales}. Evalúe esta afirmación
considerando el dato de PBI de 2023 y el papel de las expectativas.
\lineasrespuesta{10}
\begin{solucion}
Respuesta equilibrada: es cierto que parte de la desinflación fue
importada y que la contracción de 2023 ($-0{,}6\,\%$, agravada por El
Niño y conflictividad) refleja en parte el apretón monetario. Pero la
crítica omite el riesgo de desanclaje: con expectativas en 4,3\,\%, no
subir la tasa habría validado una inflación persistente y exigido un
ajuste posterior más costoso (desinflación à la Volcker). Se valora que
el estudiante discuta el dilema producto-inflación y la credibilidad
como activo del banco central.
\end{solucion}
\end{pregunta}

\begin{pregunta}[5][Propuesta]
Suponga que en 2026 ocurre un nuevo choque de oferta de magnitud similar.
Formule una recomendación de \emph{policy mix} (monetaria y fiscal) para
el Perú, indicando instrumento, oportunidad y un riesgo principal de su
propuesta.
\lineasrespuesta{10}
\begin{solucion}
Se acepta cualquier propuesta internamente consistente. Nivel excelente:
respuesta monetaria condicionada a expectativas (subir solo si se
desanclan), transferencias fiscales focalizadas y temporales a hogares
vulnerables en lugar de subsidios generalizados a combustibles
(regresivos y fiscalmente costosos), y uso de la posición fiscal
contracíclica; riesgo señalado: dominancia fiscal o persistencia del
choque que obligue a recalibrar.
\end{solucion}
\end{pregunta}

\findelexamen
\end{document}
